\chapter{Was das Modell kann und was nicht}\label{kap:grenzen}

Wahrscheinlichkeiten sind keine Vorhersagen. Wenn das Modell einer Mannschaft
30~Prozent für den Abstieg gibt und sie bleibt drin, hat es sich nicht
geirrt. Ob ein Modell gut ist, zeigt sich erst über viele Prognosen: Von
allen Ereignissen mit 30~Prozent sollten etwa drei von zehn eintreten. Dieses
Kapitel sammelt die Grenzen, die das Projekt kennt und dokumentiert hat, und
die Vereinfachungen, die sich aus dem Code ergeben.

\section{Dokumentierte Schwächen}

\subsection{Zu wenige Tore bei den Herren}

Die erwartete Torzahl je Spiel ist im Tormodell konstant $2c$. Für die
Herren sind das 2{,}64 Tore, gemessen wurden in der Bundesliga 3{,}18, rund
ein Fünftel mehr. Eine Korrektur müsste alle acht Herren-Ligen gemeinsam
treffen (Abschnitt~\ref{sec:wechselgemeinschaft}). Sie ist für das Projekt
\enquote{Prognosequalität} im Sommer 2027 vorgesehen
(\datei{docs/reports/2026-09-05-liga-empirie-zehn-ligen.md}).

\subsection{Remis: die Form der Verteilung stimmt nicht überall}\label{sec:remis}

Zwei unabhängige Poisson-Verteilungen können eine bestimmte Remisquote nicht
überschreiten. Am höchsten ist sie bei gleich starken Mannschaften ohne
Heimvorteil:
\begin{equation}
  P(\text{Remis}) \le \sum_{j \ge 0} \pi_c(j)^2 =
  \begin{cases} 26{,}1\,\% & \text{Herren } (c = 1{,}322)\\ 23{,}0\,\% & \text{Frauen } (c_F = 1{,}653).\end{cases}
\end{equation}

\begin{itemize}
  \item In der \textbf{2.~Bundesliga} lag die beobachtete Remisquote bei
        26{,}1~\%, praktisch auf dieser Obergrenze. Bei realer Streuung der
        Spielstärken kann das Modell sie nicht erreichen. Es erzeugt dort
        \emph{zu wenige} Remis.
  \item Bei den \textbf{Frauen} liegt das Modell dagegen \emph{zu hoch}:
        19{,}9~\% gegenüber beobachteten 16{,}2~\%, also 3{,}7 Prozentpunkte
        zu viel.
\end{itemize}

Die gemeinsame Ursache ist die Form der Verteilung. Die Gesamttore streuen
stärker, als Poisson erlaubt (Überdispersion 1{,}20 statt 1{,}00). Die
Verteilung der Tordifferenz ist in der Mitte dünner und an den Rändern
dicker als im Modell. Neue Parameterwerte helfen dagegen nicht. Nötig wäre
ein korreliertes Tormodell wie das von Dixon und Coles.


\subsection{Die Streuung der ELO ist kein freier Parameter}

Bei $k = 20$ und rund 34 Spielen je Saison pendelt sich die
Standardabweichung der ELO innerhalb einer Liga bei etwa 95 bis 100 ein. Wo
die echten Stärkeunterschiede größer sind, etwa in der Frauen-Bundesliga,
bildet das Modell sie zu schwach ab. Rechnerisch wären dort 458 nötig, der
Walk erreicht 240 bis 257 (ADR~0003).

\subsection{Annahmen der Regionalliga-Kopplung}

Die Kopplung beruht auf dokumentierten Annahmen
(\datei{docs/modellannahmen-rl-kopplung.md}, §\,5):

\begin{itemize}
  \item Die Zahl der Absteiger aus Nordost ist bei 2 gedeckelt.
  \item Die Nord-Meisterkopplung ist eine Näherung.
  \item West ist entkoppelt.
  \item Die Bayern-Relegation wird nicht aufgelöst.
  \item Die Verlängerung wird mit einem Drittel der Torerwartung gerechnet.
        Empirisch fallen in einer Verlängerung eher 25 bis 30~Prozent der
        Tore eines regulären Spiels.
  \item Sonderfälle wie Lizenzverweigerung, Verzicht oder Zwangsabstieg
        sind ausgelassen.
\end{itemize}

\subsection{Abstieg eines Zweitligisten und Aufstieg seiner Reserve}

Der gleichzeitige Abstieg eines Zweitligisten in die 3.~Liga und Aufstieg
seiner Zweitvertretung aus der Regionalliga wird bisher nicht
berücksichtigt. Nach §\,55b Nr.~3.1 der DFB-Spielordnung entfiele das
Aufstiegsrecht der Zweitvertretung in diesem Fall, weil ihr Verein im
kommenden Spieljahr schon mit einer Mannschaft in der 3.~Liga spielt. Dann
wäre der Nächstplatzierte ihrer Staffel aufstiegsberechtigt (ADR~0008).

Das Modell rechnet die Regionalligen ohne Malus-Lauf
(Abschnitt~\ref{sec:malus}). Um den Fall abzubilden, müsste es die
Simulation der 2.~Bundesliga mit der der Regionalliga koppeln. Die
Größenordnung ist klein: Nach einer Abschätzung vom 17.~Oktober 2026 liegt
die Wahrscheinlichkeit bei rund 1~Prozent, und praktisch betrifft sie nur
Hannover~96~II in der Regionalliga Nord.

\subsection{Keine Prognose-Historie}

Gespeichert wird nichts außer den Startwerten, und Monte Carlo ist nicht
deterministisch. Wie sich etwa die Meisterwahrscheinlichkeit im Lauf der
Saison entwickelt hat, lässt sich deshalb nicht nachträglich rekonstruieren
(ADR~0002). Der ELO-Verlauf dagegen ist deterministisch und wird bei jedem
Rendern neu berechnet.

\section{Vereinfachungen, die sich aus dem Code ergeben}

Die folgenden Punkte stehen in keinem Dokument als Schwäche, ergeben sich
aber beim Lesen des Codes. Einige sind bewusste Vereinfachungen, andere
könnten Kandidaten für Issues sein.

\begin{description}[font=\sffamily\bfseries\small,style=nextline]
  \item[Kein direkter Vergleich]
        Gleichstand nach Punkten, Tordifferenz und Toren entscheidet in der
        Simulation das alphabetisch frühere Kürzel (Abschnitt~\ref{sec:tiebreak}).
  \item[Keine Punktabzüge]
        Reale Abzüge nach Sportgerichtsurteilen oder Lizenzentscheidungen
        gehen nicht ein.
  \item[Abgesagte Spiele werden simuliert]
        Ein Spiel mit Status CANC wird in der Prognose so behandelt, als
        fände es noch statt. Das ist Absicht: Bei api-football ist CANC
        nicht verlässlich endgültig. Todesfelde gegen Norderstedt
        (Regionalliga Nord 2024/25) wurde am 21.04.2025 nachgeholt, steht
        in der API aber bis heute als CANC, ohne Ersatz-Spiel. Die
        ELO-Verlaufsseite zählt CANC nur bei der Frage, ob die Saison
        vorbei ist, als erledigt (Issue~\#292).
  \item[Rückzug eines Vereins]
        Zieht sich ein Verein während der Saison zurück (Türkgücü München,
        3.~Liga 2021/22; Türkspor Dortmund und KFC Uerdingen, Regionalliga
        West 2024/25), setzt die API seine Restspiele auf CANC. Seine
        bereits gespielten Spiele bleiben aber als beendet stehen, obwohl
        der Verband sie annulliert. Auch ein kampflos gewertetes Spiel kann
        als CANC statt als AWD kommen. Die Prognose bildet einen Rückzug
        deshalb nicht richtig ab; dafür braucht es einen Eingriff von Hand
        (Issue~\#292).
  \item[Kein Rückgang zur Mitte]
        Die ELO wird beim Saisonwechsel 1:1 übertragen. Viele ELO-Systeme
        ziehen die Werte zwischen zwei Saisons ein Stück zum Mittel zurück,
        um Kaderwechsel abzubilden.
  \item[Zwei Erwartungsskalen]
        ELO-Update und Tormodell rechnen mit unterschiedlich skalierten
        Erwartungen (Abschnitt~\ref{sec:skalen}).
  \item[Gewertete Spiele in der Playoff-ELO]
        Die aktuelle ELO für die Aufstiegsspiele
        (\code{rl\_aktuelle\_elo}, \datei{RCode/rl_verdrahtung.R:113--121})
        wird ohne das Kennzeichen \code{elo\_neutral} berechnet. Ein am grünen
        Tisch gewertetes Spiel würde dort die ELO bewegen, anders als
        überall sonst. Ob das beabsichtigt ist, ist im Code nicht kommentiert.
\end{description}
